\scanpage{351}
% top centre
\mspage{267}
% lines 1–15
conjecture que sous ces conditions, $\underline{\unsure{\Gamma|\mathscr T}}(N/J)\overset{\mathrm{def}}{=}\underline{\mathrm{Hom}}_J(J,N)$ est accessible, et qu'il en est de $\unsure{\overline{\mathrm m}|\overline{\mathrm w}}$ de $\underline{\mathrm{Hom}}\,\underline{\unsure{\mathrm{cart}|\mathrm{cat}}}_J(J,N)$. C'est ce qui est prouvé dans SGA 4 I 9.2\deleted{\unsure{2|7}}\unsure{4|7}. sous des hypothèses légèrement plus fortes (à l'aide de l'ancienne notion, trop forte à mon gré, de filtration cardinale). Je serais étonné que la démonstration (un peu \unsure{longuette}) de loc. cit. ne s'applique pas au cas général\unsure{.|?} Mais on aimerait avoir un énoncé qui \unsure{coiffe} celui-ci et 4.10.8, dans le cas d'un foncteur $N\longrightarrow J$ qui \deleted{n'est}\inserted{ne serait} supposé ni fibrant ni cofibrant, mais \unsure{....|...} (???).

% lines 16–17, ruled heading
\noindent\vrule\hspace{0.5em}\parbox{0.93\linewidth}{4.11. \unsure{Application|Applications} à l'existence des limites dans une catégorie de la forme $\operatorname{Ind}_{\pi_0}(\unsure{C|\mathcal C})$.}

% lines 18–19
\uline{Proposition} 4.11.1. \deleted{Soit $\unsure{\Phi|\underline{\Phi}}$ un ensemble (pas \unsure{néc} petit) de petites catégories. \unsure{On}} Soit
% archive sheet number: 350
